% ======================================================================
% BLUEPRINT FOR THE PROOF OF PROPOSITION~\ref{prop:cavity}
% ======================================================================
%
% PURPOSE.
% Complete the proof of Proposition~\ref{prop:cavity} rigorously, filling in
% all cavity-interpolation computations and all error estimates.  The final
% proof should be self-contained at the level of the paper: standard Gaussian
% integration by parts, Holder, Young, Taylor with integral remainder, and
% Gronwall may be invoked, but every nontrivial application must be written
% down sufficiently explicitly that the constants and powers of N can be
% checked.
%
% ----------------------------------------------------------------------
% FILES IN research_public TO CONSULT
% ----------------------------------------------------------------------
%
% The main paper source is
%
%   latalaStrictAlmeidaThouless/refs/
%       latalaArgumentsStrictAlmeidaThoulessCondition.tex
%
% This is the authoritative source for the notation and for the statement of
% Proposition~\ref{prop:cavity}.  Preserve its notation and equation labels.
% In particular, do not introduce a second convention for A_s,B_s,C_s or for
% U_s,V_s,D_s.
%
% The most important consistency check is
%
%   latalaStrictAlmeidaThouless/Lemmas/ATDefs.lean
%
% Search there for
%
%   cavityKappa
%   cavityZeta
%   theta
%   cavityMatrix
%   cavityChangeMatrix
%   cavityRemainder
%   cavityErrorScale
%
% The definitions in that file encode exactly the coefficient calculation
% which must emerge from the last-spin interpolation.  In particular, with
%
%   b_2 = beta^2(1-q^2),
%   b_1 = beta^2(q-q^2),
%   b_0 = beta^2(r-q^2),
%
% the matrix in the (A,B,C)-coordinates is
%
%   [ b_2,             -4b_1,               3b_0              ]
%   [ b_1, b_2-2b_1-3b_0,               6b_0-3b_1             ]
%   [ b_0,       4b_1-8b_0, b_2-8b_1+10b_0                   ].
%
% Use this as a check on the replica combinatorics below.  Do not merely cite
% the Lean file in the paper proof: reproduce the mathematical calculation.
%
% Also consult
%
%   latalaStrictAlmeidaThouless/Lemmas/Cavity/Talagrand_Cavity.lean
%
% Search for
%
%   cavityChangeMatrix_mul_cavityMatrix
%   cavityChangeMatrix_mulVec_theta
%   cavityModeMatrix
%   cavityModeRemainder_bound_from_lastSpin
%   beta_sq_mul_repliconCoefficient_eq_atParameter
%
% The deterministic mode algebra after the last-spin calculation is already
% checked there.  In particular, note carefully that the Lean file uses the
% row order (V,U,D), not (U,V,D).  It verifies that the cavity matrix becomes
% triangular:
%
%   V -> beta^2(kappa V + zeta U),
%   U -> beta^2 kappa U,
%   D -> beta^2(1-2q+r)D.
%
% Therefore, if a hand calculation below produces any different coefficient,
% regard that as a likely replica-counting error and debug it before proceeding.
%
% IMPORTANT: do not use
%
%   cavityModeRemainder_bound_from_lastSpin
%
% itself as a black box.  That theorem is precisely the analytic statement
% whose mathematical proof is being supplied here.  Use only the definitions
% and deterministic algebra surrounding it as consistency checks.
%
% Consult
%
%   latalaStrictAlmeidaThouless/Lemmas/fixed_point.lean
%
% for the canonical identities concerning q,r and the AT parameter.  In
% particular, use as a consistency check
%
%   alpha = beta^2(1-2q+r).
%
% The file also records q as the RS fixed point and r as the fourth tanh
% moment.  In the paper proof these identities should continue to be derived
% from the already stated definitions rather than from Lean.
%
% The file
%
%   latalaStrictAlmeidaThouless/Lemmas/README.md
%
% confirms the intended dependency order: Gaussian differentiation first,
% then the cavity coefficients/stability, and then the remaining argument.
% It explicitly says that the last-spin interpolation is one of the analytic
% constructions still requiring proof.
%
% The file
%
%   latalaStrictAlmeidaThouless/Lemmas/Cavity/cavity_interpolation.lean
%
% is empty at present.  Do not spend time looking there for a proof.
%
% Avoid using the final overlap concentration theorem, the final absorption
% argument, or the conclusion of Theorem~\ref{thm:main} anywhere in this
% proof.  Proposition~\ref{prop:cavity} must precede those arguments.
%
% In particular, the desired error is intentionally left in the form
%
%   N^{-3/2} + nu_s[|Q_{12}|^3].
%
% Do not replace the cubic moment by a later concentration estimate inside
% this proposition.
%
% ======================================================================
% START OF THE ACTUAL PROOF
% ======================================================================

\begin{proof}[Proof of Proposition~\ref{prop:cavity}]
All implicit constants below are uniform for
$(\beta,h,s)\in\cK\times[0,1]$ and depend only on $\cK$.
For convenience set
\[
 \eta_{N,s}
 \coloneqq
 N^{-3/2}+\nu_s[|Q_{12}|^3].
\]
Since $\cK$ is compact, $\beta$ is uniformly bounded on $\cK$.
Moreover $q,r\in[0,1]$, so every polynomial expression in
$\beta,q,r$ occurring below is uniformly bounded.

\paragraph{Exchangeability and removal of the last site.}

Put
\[
 \varepsilon_a\coloneqq\sigma_N^a,
 \qquad
 Q^-_{ab}
 \coloneqq
 \frac1N\sum_{i<N}\sigma_i^a\sigma_i^b-q.
\]
Then
\begin{equation}
 Q_{ab}
 =
 Q^-_{ab}+\frac1N\varepsilon_a\varepsilon_b.
 \label{eq:Qab}
\end{equation}

First justify carefully the three site-exchangeability identities used below.
Since
\[
 Q_{ab}
 =
 \frac1N\sum_{i=1}^N
 \bigl(\sigma_i^a\sigma_i^b-q\bigr),
\]
we have
\[
 A_s
 =
 \frac1N\sum_{i=1}^N
 \nu_s\!\left[
 Q_{12}(\sigma_i^1\sigma_i^2-q)
 \right].
\]
After the disorder average, the joint law is invariant under permutations of
the sites, so every summand has the same value.  Therefore
\[
 A_s
 =
 \nu_s\!\left[
 Q_{12}(\varepsilon_1\varepsilon_2-q)
 \right].
\]
Exactly the same argument gives
\begin{align}
 A_s&=\nu_s\!\left[Q_{12}(\varepsilon_1\varepsilon_2-q)\right],
 \nonumber\\
 B_s&=\nu_s\!\left[Q_{12}(\varepsilon_1\varepsilon_3-q)\right],
 \label{eq:cavity-site-exchangeability}\\
 C_s&=\nu_s\!\left[Q_{12}(\varepsilon_3\varepsilon_4-q)\right].
 \nonumber
\end{align}

\paragraph{Definition of the last-spin interpolation.}

Write $\sigma=(\rho,\varepsilon)$ with
$\rho\in\{-1,1\}^{N-1}$ and $\varepsilon\in\{-1,1\}$.
Only the random interaction involving the last site is interpolated.
Replace it by
\[
 \frac{\beta\sqrt{su}}{\sqrt N}
 \sum_{i<N}
 \frac{g_{iN}+g_{Ni}}{\sqrt2}\rho_i\varepsilon
 +
 \beta\sqrt{s(1-u)q}\,\widehat z\,\varepsilon,
 \qquad 0\le u\le1,
\]
where $\widehat z$ is independent of all previously introduced disorder.

Denote the associated disorder-averaged replicated Gibbs expectation by
$\nu_{s,u}$.  Check directly from the Hamiltonian that
\[
 \nu_{s,1}=\nu_s.
\]

At $u=0$, all random variables entering the last-spin field are independent
of all disorder variables entering the Hamiltonian of the first $N-1$ spins.
The last-spin field equals
\[
 h+\beta\sqrt{(1-s)q}\,z_N+\beta\sqrt{sq}\,\widehat z.
\]
Since $z_N$ and $\widehat z$ are independent standard Gaussian variables,
\[
 \sqrt{1-s}\,z_N+\sqrt{s}\,\widehat z
 \stackrel{\mathrm d}=Z,
 \qquad Z\sim N(0,1),
\]
and hence the last-spin field has distribution
\[
 h+\beta\sqrt q\,Z.
\]

Conditionally on this field, distinct replicas of the last spin are
independent and have mean
\[
 \tanh(h+\beta\sqrt q\,Z).
\]
Consequently, using the definitions of $q$ and $r$,
\begin{equation}
 \mathbb E_0[\varepsilon_a\varepsilon_b]=q,
 \qquad
 \mathbb E_0[
 \varepsilon_a\varepsilon_b\varepsilon_c\varepsilon_d
 ]=r
 \label{eq:last-spin-moments-proof}
\end{equation}
whenever the displayed replica indices are distinct.

Make explicit here that the last-spin variables are also independent of
observables depending only on the first $N-1$ spins, after averaging over the
independent last-spin disorder.

\paragraph{Gaussian differentiation formula.}

For configurations $(\rho,\varepsilon)$ and
$(\rho',\varepsilon')$, compute the covariance of the interpolated random
last-spin field.  The derivative with respect to $u$ is
\[
 s\beta^2\varepsilon\varepsilon'
 \left(
 \frac1N\sum_{i<N}\rho_i\rho_i'-q
 \right)
 =
 s\beta^2\varepsilon\varepsilon'Q^-(\rho,\rho').
\]
For coincident configurations its value is
\[
 s\beta^2
 \left(\frac{N-1}{N}-q\right),
\]
which does not depend on the spin configuration.  Thus the diagonal
covariance derivative disappears from the normalized Gibbs derivative.

Apply Gaussian integration by parts to the numerator and denominator of the
$n$-replica Gibbs expectation.  Write the full normalized calculation, or
quote the standard Gaussian Gibbs differentiation identity and substitute
the covariance derivative above.  The result must be
\begin{align}
 \frac{\mathrm d}{\mathrm du}\nu_{s,u}[F]
 =s\beta^2\Bigg\{&
 \sum_{1\le a<b\le n}
 \nu_{s,u}
 [F\varepsilon_a\varepsilon_bQ^-_{ab}]
 \nonumber\\
 &-n\sum_{a=1}^n
 \nu_{s,u}
 [F\varepsilon_a\varepsilon_{n+1}Q^-_{a,n+1}]
 \nonumber\\
 &+\frac{n(n+1)}2
 \nu_{s,u}
 [F\varepsilon_{n+1}\varepsilon_{n+2}
 Q^-_{n+1,n+2}]
 \Bigg\}.
 \label{eq:cavity-derivative-proof}
\end{align}

Codex should check the factors
\[
 1,\qquad -n,\qquad \frac{n(n+1)}2
\]
rather than silently importing them from the draft.

Replica labels may always be enlarged when necessary: if $F$ originally
depends on fewer replicas, regard it as a function of a larger replica family
which ignores the additional coordinates.

\paragraph{Uniform cubic cavity-overlap estimate.}

Set
\[
 M_3(u)
 \coloneqq
 \nu_{s,u}[|Q^-_{12}|^3].
\]
Use \eqref{eq:cavity-derivative-proof} with
\[
 F=|Q^-_{12}|^3.
\]
Since
\[
 |Q^-_{ab}|
 \le
 \frac{N-1}{N}+q
 \le2,
\]
every term in the differentiated expression is bounded by a fixed constant
times $M_3(u)$.  For example,
\[
 \left|
 \nu_{s,u}\!
 \left[
 |Q^-_{12}|^3
 \varepsilon_a\varepsilon_bQ^-_{ab}
 \right]
 \right|
 \le
 2M_3(u).
\]
There are only finitely many terms, and $\beta$ is uniformly bounded on
$\cK$, hence
\begin{equation}
 |M_3'(u)|\le C_{\cK}M_3(u).
 \label{eq:M3-differential}
\end{equation}

At $u=1$, \eqref{eq:Qab} and
\[
 |x+y|^3\le4(|x|^3+|y|^3)
\]
give
\[
 |Q^-_{12}|^3
 \le
 4|Q_{12}|^3+\frac4{N^3}.
\]
Thus
\[
 M_3(1)
 \le
 4\nu_s[|Q_{12}|^3]+\frac4{N^3}.
\]

Apply Gronwall backward from $u=1$.  More explicitly,
\eqref{eq:M3-differential} implies
\[
 M_3(u)
 \le
 e^{C_{\cK}(1-u)}M_3(1),
 \qquad 0\le u\le1.
\]
Since $N^{-3}\le N^{-3/2}$ for $N\ge1$,
\begin{equation}
 \sup_{0\le u\le1}M_3(u)
 \le
 C_{\cK}\eta_{N,s}.
 \label{eq:M3-gronwall-proof}
\end{equation}

\paragraph{Replacement of quadratic overlap products by their $u=0$
cavity versions.}

Let
\[
 P^-=Q^-_{ab}Q^-_{cd},
 \qquad
 (ab,cd)\in
 \{(12,12),(12,13),(12,34)\}.
\]
Apply \eqref{eq:cavity-derivative-proof} to $P^-$.
Every differentiated term is, up to a bounded spin factor, a product of
three cavity overlaps:
\[
 Q^-_{ab}Q^-_{cd}Q^-_{ef}.
\]
By Holder and replica exchangeability,
\begin{align*}
 \nu_{s,u}
 \left[
 |Q^-_{ab}Q^-_{cd}Q^-_{ef}|
 \right]
 &\le
 \nu_{s,u}[|Q^-_{ab}|^3]^{1/3}
 \nu_{s,u}[|Q^-_{cd}|^3]^{1/3}
 \nu_{s,u}[|Q^-_{ef}|^3]^{1/3}\\
 &=M_3(u).
\end{align*}
The same estimate remains valid when two of the replica edges coincide.
Hence
\[
 \left|
 \frac{\mathrm d}{\mathrm du}\nu_{s,u}[P^-]
 \right|
 \le
 C_{\cK}M_3(u).
\]
Integrating and using \eqref{eq:M3-gronwall-proof},
\begin{equation}
 \left|
 \nu_{s,0}[P^-]-\nu_{s,1}[P^-]
 \right|
 \le C_{\cK}\eta_{N,s}.
 \label{eq:Pminus-endpoints}
\end{equation}

It remains to replace $Q^-Q^-$ at $u=1$ by $QQ$.
From
\[
 Q^-_{ab}
 =
 Q_{ab}-\frac1N\varepsilon_a\varepsilon_b
\]
one gets
\[
 \left|
 Q^-_{ab}Q^-_{cd}-Q_{ab}Q_{cd}
 \right|
 \le
 \frac{|Q_{ab}|+|Q_{cd}|}{N}
 +\frac1{N^2}.
\]
Use Young's inequality with exponents $3$ and $3/2$:
\begin{equation}
 \frac{|x|}{N}
 \le
 \frac{|x|^3}{3}
 +\frac{2}{3N^{3/2}}.
 \label{eq:young-cavity}
\end{equation}
Replica exchangeability and $N^{-2}\le N^{-3/2}$ then imply
\begin{equation}
 \left|
 \nu_{s,0}[P^-]-\nu_s[Q_{ab}Q_{cd}]
 \right|
 \le C_{\cK}\eta_{N,s}.
 \label{eq:endpoint-replacement-proof}
\end{equation}

Define
\begin{align*}
 A_s^-&\coloneqq\nu_{s,0}[(Q^-_{12})^2],\\
 B_s^-&\coloneqq\nu_{s,0}[Q^-_{12}Q^-_{13}],\\
 C_s^-&\coloneqq\nu_{s,0}[Q^-_{12}Q^-_{34}],
\end{align*}
and
\begin{align*}
 U_s^-&\coloneqq A_s^--4B_s^-+3C_s^-,\\
 V_s^-&\coloneqq2B_s^--3C_s^-,\\
 D_s^-&\coloneqq A_s^--2B_s^-+C_s^-.
\end{align*}
Since these are fixed linear combinations,
\begin{equation}
 |U_s^--U_s|
 +|V_s^--V_s|
 +|D_s^--D_s|
 \le
 C_{\cK}\eta_{N,s}.
 \label{eq:cavity-mode-replacement}
\end{equation}

\paragraph{Separation into off-diagonal and diagonal last-spin terms.}

Using \eqref{eq:Qab} in
\eqref{eq:cavity-site-exchangeability}, write
\begin{align}
 A_s&=\nu_s[G_A]+\frac1N\nu_s[E_A],
 \nonumber\\
 B_s&=\nu_s[G_B]+\frac1N\nu_s[E_B],
 \label{eq:cavity-decomposition-proof}\\
 C_s&=\nu_s[G_C]+\frac1N\nu_s[E_C],
 \nonumber
\end{align}
where
\begin{align*}
 G_A&=Q^-_{12}(\varepsilon_1\varepsilon_2-q),
 &
 E_A&=\varepsilon_1\varepsilon_2
      (\varepsilon_1\varepsilon_2-q),
 \\
 G_B&=Q^-_{12}(\varepsilon_1\varepsilon_3-q),
 &
 E_B&=\varepsilon_1\varepsilon_2
      (\varepsilon_1\varepsilon_3-q),
 \\
 G_C&=Q^-_{12}(\varepsilon_3\varepsilon_4-q),
 &
 E_C&=\varepsilon_1\varepsilon_2
      (\varepsilon_3\varepsilon_4-q).
\end{align*}

Put
\[
 g_j(u)=\nu_{s,u}[G_j],
 \qquad
 e_j(u)=\nu_{s,u}[E_j],
 \qquad j\in\{A,B,C\}.
\]
Use
\begin{align*}
 g_U&=g_A-4g_B+3g_C,
 &
 e_U&=e_A-4e_B+3e_C,
 \\
 g_V&=2g_B-3g_C,
 &
 e_V&=2e_B-3e_C,
 \\
 g_D&=g_A-2g_B+g_C,
 &
 e_D&=e_A-2e_B+e_C.
\end{align*}

At $u=0$, $Q^-_{12}$ depends only on the first $N-1$ spins whereas the
centered last-spin factor is independent of those spins and has expectation
zero:
\[
 \mathbb E_0[
 \varepsilon_p\varepsilon_q-q
 ]=0.
\]
Therefore
\begin{equation}
 g_A(0)=g_B(0)=g_C(0)=0.
 \label{eq:g-zero}
\end{equation}

\paragraph{Second derivative bound and Taylor expansion.}

Applying \eqref{eq:cavity-derivative-proof} once to a $G_j$ introduces
one additional cavity overlap.  Applying it a second time introduces one more.
Thus $g_j''(u)$ is a finite linear combination of terms of the form
\[
 \nu_{s,u}
 \left[
 \Xi(\varepsilon)
 Q^-_{a_1b_1}Q^-_{a_2b_2}Q^-_{a_3b_3}
 \right],
\]
where $|\Xi|\le C$ and the number of replicas involved is bounded by an
absolute constant.

By Holder and replica exchangeability,
\[
 |g_j''(u)|
 \le C_{\cK}M_3(u)
 \le C_{\cK}\eta_{N,s}.
\]
The same estimate holds for $g_U'',g_V'',g_D''$.

Use Taylor's formula with integral remainder:
\[
 g_X(1)
 =
 g_X(0)+g_X'(0)
 +\int_0^1(1-u)g_X''(u)\,\mathrm du.
\]
Together with \eqref{eq:g-zero}, this yields
\begin{equation}
 g_X(1)
 =
 g_X'(0)+O_{\cK}(\eta_{N,s}),
 \qquad
 X\in\{U,V,D\}.
 \label{eq:cavity-Taylor-modes}
\end{equation}

Do not simply write ``applying the derivative identity twice'' without this
argument.  The proof must explain why exactly three cavity overlaps occur and
why Holder reduces every such term to $M_3(u)$.

\paragraph{The three-valued last-spin edge rule.}

Introduce
\[
 t_2\coloneqq1-q^2,
 \qquad
 t_1\coloneqq q-q^2,
 \qquad
 t_0\coloneqq r-q^2.
\]
For distinct $p,q$ and a replica edge $\{a,b\}$,
\eqref{eq:last-spin-moments-proof} gives
\begin{equation}
 \mathbb E_0[
 (\varepsilon_p\varepsilon_q-q)
 \varepsilon_a\varepsilon_b
 ]
 =
 \begin{cases}
 t_2,
 &\{a,b\}=\{p,q\},\\
 t_1,
 &|\{a,b\}\cap\{p,q\}|=1,\\
 t_0,
 &\{a,b\}\cap\{p,q\}=\varnothing.
 \end{cases}
 \label{eq:edge-rule-proof}
\end{equation}

Spell out the middle case, since it is useful for checking signs:
if, for example, $a=p$ and $b\notin\{p,q\}$, then
\[
 (\varepsilon_p\varepsilon_q-q)
 \varepsilon_p\varepsilon_b
 =
 \varepsilon_q\varepsilon_b
 -q\varepsilon_p\varepsilon_b,
\]
whose expectation equals
\[
 q-q^2.
\]

\paragraph{Explicit replica counting for $g_A'(0),g_B'(0),g_C'(0)$.}

This is the main combinatorial calculation and should not be suppressed.

Regard every $G_j$ as a function of four replicas and apply
\eqref{eq:cavity-derivative-proof} with $n=4$.  The derivative contains the
edges
\[
 \{a,b\}\subset\{1,2,3,4\}
 \quad\text{with coefficient }+1,
\]
the edges
\[
 \{a,5\},
 \quad 1\le a\le4,
 \quad\text{with coefficient }-4,
\]
and
\[
 \{5,6\}
 \quad\text{with coefficient }10.
\]

For the overlap part, classify the new edge relative to $\{1,2\}$:
\[
 \begin{array}{c|c}
 \text{relation to }\{1,2\}
 &
 \nu_{s,0}[Q^-_{12}Q^-_{ab}]
 \\ \hline
 \{a,b\}=\{1,2\}
 &
 A_s^-
 \\
 |\{a,b\}\cap\{1,2\}|=1
 &
 B_s^-
 \\
 \{a,b\}\cap\{1,2\}=\varnothing
 &
 C_s^- .
 \end{array}
\]
The equalities follow from replica exchangeability.

For the spin part, independently classify the same edge relative to the
edge appearing in $G_j$ and use
\eqref{eq:edge-rule-proof}.

The resulting calculation must give
\begin{align}
 g_A'(0)
 &=
 s\beta^2
 \bigl(
 t_2 A_s^- -4t_1B_s^-+3t_0C_s^-
 \bigr),
 \label{eq:gA-prime-explicit}\\
 g_B'(0)
 &=
 s\beta^2
 \bigl(
 t_1A_s^-
 +(t_2-2t_1-3t_0)B_s^-
 +(6t_0-3t_1)C_s^-
 \bigr),
 \label{eq:gB-prime-explicit}\\
 g_C'(0)
 &=
 s\beta^2
 \bigl(
 t_0A_s^-
 +(4t_1-8t_0)B_s^-
 +(t_2-8t_1+10t_0)C_s^-
 \bigr).
 \label{eq:gC-prime-explicit}
\end{align}

These three formulas are exactly the matrix calculation encoded by
\texttt{cavityMatrix} in
\texttt{Lemmas/ATDefs.lean}.  Use that file to check every coefficient.

A particularly transparent way to verify
\eqref{eq:gA-prime-explicit} is as follows.
For $G_A$, the centered spin edge is $\{1,2\}$.  Among the derivative edges,
the contribution associated with an $A_s^-$ overlap has coefficient $t_2$,
the total contribution associated with a $B_s^-$ overlap is $-4t_1$, and the
total contribution associated with a $C_s^-$ overlap is $3t_0$.
Perform the analogous edge count for $G_B$ and $G_C$ rather than guessing
their coefficients.

\paragraph{Pass from $(A,B,C)$ to the three cavity modes.}

Recall
\[
 \kappa=1-4q+3r,
 \qquad
 \zeta=2q+q^2-3r.
\]
The elementary identities
\begin{align}
 t_2-4t_1+3t_0&=\kappa,
 \label{eq:kappa-t}\\
 2t_1-3t_0&=\zeta,
 \label{eq:zeta-t}\\
 t_2-2t_1+t_0&=1-2q+r
 \label{eq:replicon-t}
\end{align}
should be displayed.

Now form
\[
 g_U'=g_A'-4g_B'+3g_C',
 \qquad
 g_V'=2g_B'-3g_C',
 \qquad
 g_D'=g_A'-2g_B'+g_C'.
\]
A direct substitution of
\eqref{eq:gA-prime-explicit}--\eqref{eq:gC-prime-explicit}
and
\eqref{eq:kappa-t}--\eqref{eq:replicon-t}
gives
\begin{align}
 g_U'(0)
 &=
 s\beta^2\kappa U_s^-,
 \nonumber\\
 g_V'(0)
 &=
 s\beta^2
 \bigl(
 \zeta U_s^-+\kappa V_s^-
 \bigr),
 \label{eq:scalar-first-derivatives}\\
 g_D'(0)
 &=
 s\beta^2(1-2q+r)D_s^-
 =
 s\alpha D_s^-.
 \nonumber
\end{align}

For checking this calculation against the repository, use
\texttt{cavityChangeMatrix\_mul\_cavityMatrix} in
\texttt{Lemmas/Cavity/Talagrand\_Cavity.lean}.
Remember that the Lean file orders the transformed vector as
\[
 (V,U,D),
\]
so its triangular matrix reads
\[
 \begin{pmatrix}
 \beta^2\kappa&\beta^2\zeta&0\\
 0&\beta^2\kappa&0\\
 0&0&\beta^2(1-2q+r)
 \end{pmatrix}.
\]
This is equivalent to \eqref{eq:scalar-first-derivatives}.

\paragraph{The diagonal $N^{-1}$ terms.}

At $u=0$,
\begin{align*}
 e_A(0)
 &=
 \mathbb E_0[
 \varepsilon_1\varepsilon_2
 (\varepsilon_1\varepsilon_2-q)
 ]
 =
 1-q^2
 =
 t_2,
 \\
 e_B(0)
 &=
 \mathbb E_0[
 \varepsilon_1\varepsilon_2
 (\varepsilon_1\varepsilon_3-q)
 ]
 =
 q-q^2
 =
 t_1,
 \\
 e_C(0)
 &=
 \mathbb E_0[
 \varepsilon_1\varepsilon_2
 (\varepsilon_3\varepsilon_4-q)
 ]
 =
 r-q^2
 =
 t_0.
\end{align*}
Therefore
\begin{align}
 e_U(0)
 &=
 t_2-4t_1+3t_0
 =
 \kappa,
 \nonumber\\
 e_V(0)
 &=
 2t_1-3t_0
 =
 \zeta,
 \label{eq:scalar-diagonal-terms}\\
 e_D(0)
 &=
 t_2-2t_1+t_0
 =
 1-2q+r
 =
 \frac{\alpha}{\beta^2}.
 \nonumber
\end{align}

This calculation is the scalar version of
\texttt{cavityChangeMatrix\_mulVec\_theta} in
\texttt{Lemmas/Cavity/Talagrand\_Cavity.lean}, where
\[
 \theta=(t_2,t_1,t_0).
\]

\paragraph{Error made by replacing $e_X(1)$ by $e_X(0)$.}

Each $E_A,E_B,E_C$ is uniformly bounded, since
$\varepsilon_a\in\{-1,1\}$ and $0\le q\le1$.
Apply \eqref{eq:cavity-derivative-proof}.  Every derivative term contains
exactly one cavity overlap, multiplied by a bounded spin factor.  Hence,
using replica exchangeability,
\[
 |e_X'(u)|
 \le
 C_{\cK}\nu_{s,u}[|Q^-_{12}|],
 \qquad
 X\in\{U,V,D\}.
\]
Therefore
\[
 \frac1N|e_X(1)-e_X(0)|
 \le
 \frac{C_{\cK}}N
 \int_0^1
 \nu_{s,u}[|Q^-_{12}|]\,\mathrm du.
\]

Apply the pointwise Young inequality
\[
 \frac{|Q^-_{12}|}{N}
 \le
 \frac{|Q^-_{12}|^3}{3}
 +\frac{2}{3N^{3/2}}.
\]
After averaging and using \eqref{eq:M3-gronwall-proof},
\begin{equation}
 \frac1N|e_X(1)-e_X(0)|
 \le
 C_{\cK}\eta_{N,s},
 \qquad X\in\{U,V,D\}.
 \label{eq:scalar-diagonal-remainders}
\end{equation}

\paragraph{Assembly of the $U$ equation.}

From the decomposition
\eqref{eq:cavity-decomposition-proof},
\[
 U_s
 =
 g_U(1)+\frac1N e_U(1).
\]
Use successively
\eqref{eq:cavity-Taylor-modes},
\eqref{eq:scalar-first-derivatives},
\eqref{eq:scalar-diagonal-terms},
and \eqref{eq:scalar-diagonal-remainders}:
\[
 U_s
 =
 s\beta^2\kappa U_s^-
 +\frac{\kappa}{N}
 +O_{\cK}(\eta_{N,s}).
\]
Finally use \eqref{eq:cavity-mode-replacement}.  Since
$s\beta^2\kappa$ is uniformly bounded,
\[
 s\beta^2\kappa(U_s^--U_s)
 =
 O_{\cK}(\eta_{N,s}),
\]
and hence
\begin{equation}
 U_s
 =
 s\beta^2\kappa U_s
 +\frac{\kappa}{N}
 +O_{\cK}(\eta_{N,s}).
 \label{eq:assembled-U}
\end{equation}

\paragraph{Assembly of the $V$ equation.}

Likewise,
\[
 V_s
 =
 g_V(1)+\frac1N e_V(1),
\]
so
\[
 V_s
 =
 s\beta^2(\zeta U_s^-+\kappa V_s^-)
 +\frac{\zeta}{N}
 +O_{\cK}(\eta_{N,s}).
\]
By \eqref{eq:cavity-mode-replacement} and the uniform boundedness of
$\beta,\kappa,\zeta$,
\begin{equation}
 V_s
 =
 s\beta^2(\zeta U_s+\kappa V_s)
 +\frac{\zeta}{N}
 +O_{\cK}(\eta_{N,s}).
 \label{eq:assembled-V}
\end{equation}

\paragraph{Assembly of the replicon equation.}

Similarly,
\[
 D_s
 =
 g_D(1)+\frac1N e_D(1),
\]
and therefore
\[
 D_s
 =
 s\alpha D_s^-
 +\frac{\alpha}{\beta^2N}
 +O_{\cK}(\eta_{N,s}).
\]
Using \eqref{eq:cavity-mode-replacement},
\begin{equation}
 D_s
 =
 s\alpha D_s
 +\frac{\alpha}{\beta^2N}
 +O_{\cK}(\eta_{N,s}).
 \label{eq:assembled-D}
\end{equation}

\paragraph{Definition of the remainders and conclusion.}

Rearrange
\eqref{eq:assembled-U}--\eqref{eq:assembled-D} as
\begin{align*}
 (1-s\beta^2\kappa)U_s
 &=
 \frac{\kappa}{N}+\mathcal R^U_{N,s},
 \\
 (1-s\beta^2\kappa)V_s
 -s\beta^2\zeta U_s
 &=
 \frac{\zeta}{N}+\mathcal R^V_{N,s},
 \\
 (1-s\alpha)D_s
 &=
 \frac{\alpha}{\beta^2N}
 +\mathcal R^D_{N,s}.
\end{align*}
The preceding estimates give, uniformly on
$\cK\times[0,1]$,
\[
 \max_{X\in\{U,V,D\}}
 |\mathcal R^X_{N,s}|
 \le
 C_{\cK}
 \left(
 N^{-3/2}+\nu_s[|Q_{12}|^3]
 \right).
\]
These are exactly
\eqref{eq:cavity-U}--\eqref{eq:cavity-remainders}.
\end{proof}

% ======================================================================
% CODEX CHECKLIST BEFORE ACCEPTING THE COMPLETED PROOF
% ======================================================================
%
% The final proof is not complete unless all of the following are checked.
%
% 1. The covariance derivative of the last-spin interpolation is computed
%    explicitly and equals
%
%       s beta^2 epsilon epsilon' Q^-.
%
% 2. The normalized Gaussian derivative formula has coefficients
%
%       1, -n, n(n+1)/2.
%
% 3. The proof of the M_3 estimate explains why differentiation does not
%    require a fourth-moment estimate: use |Q^-| <= 2.
%
% 4. The endpoint comparison for quadratic products uses Holder for products
%    of three cavity overlaps and Young for |Q|/N.
%
% 5. The Taylor remainder for g_X is written using g_X'' and is bounded
%    solely by the cubic overlap moment.
%
% 6. The first derivatives g_A'(0),g_B'(0),g_C'(0) are checked against
%    cavityMatrix in ATDefs.lean.  In particular the three rows must be
%
%       (t_2, -4t_1, 3t_0),
%
%       (t_1, t_2-2t_1-3t_0, 6t_0-3t_1),
%
%       (t_0, 4t_1-8t_0, t_2-8t_1+10t_0).
%
% 7. After changing coordinates, the mode equations must agree with
%    Talagrand_Cavity.lean:
%
%       U -> beta^2 kappa U,
%
%       V -> beta^2(zeta U+kappa V),
%
%       D -> beta^2(1-2q+r)D.
%
% 8. The source vector must become
%
%       (kappa,zeta,1-2q+r)
%
%    in the paper's (U,V,D) order.
%
% 9. Use alpha = beta^2(1-2q+r) only after the coefficient
%    1-2q+r has been obtained from the cavity calculation.
%
% 10. Every O-term is bounded by
%
%       C_K (N^{-3/2} + nu_s[|Q_{12}|^3])
%
%     with a constant uniform in s and in (beta,h) in K.
%
% 11. Do not invoke the final overlap concentration estimate or any theorem
%     downstream of Proposition~\ref{prop:cavity}.
%
% 12. Do not cite cavityModeRemainder_bound_from_lastSpin as the proof:
%     that theorem is the formal placeholder corresponding to the analytic
%     argument being supplied here.
% ======================================================================